\documentclass[12pt,uplatex,dvipdfmx,book,oneside]{jlreq}
\usepackage{amssymb,amsmath,latexsym}
\usepackage[haranoaji]{pxchfon} % 原ノ味フォントを使う

\begin{document}
\pagenumbering{roman}

\begin{center}
{\Large
令和9年度 修士論文
}

\vspace{4cm}
{\Huge
和文題名（長さにより文字の大きさを調節して下さい）
}

\vspace{12pt}

{\huge
Title in English
}
\end{center}

\vfill
\begin{flushright}
\large
日本大学大学院理工学研究科博士前期課程数学専攻

%学生番号
M6012
%氏名
川村聡一郎

指導教員 水野 将司
\end{flushright}

\thispagestyle{empty}

\newpage

\tableofcontents
\thispagestyle{empty}


\chapter{○○}
\pagenumbering{arabic}





\section{○○○}

\subsection{界面モビリティの計算}
\begin{align}
    \phi_\mathrm{eq}
    &= \frac{1}{2}\left(1- \tanh{\frac{\sqrt{2W}}{2a}}\right)\\
    \frac{d\phi_\mathrm{eq}}{dx}
    &= \frac{1}{2}\left(1- \frac{4}{(e^t+ e^{-t})^2}\right)\ \ \ t= \frac{\sqrt{2W}}{2a}x\\
    \left(\frac{d\phi_\mathrm{eq}}{dx}\right)^2
    &= \frac{1}{4}\left(1- \frac{8}{(e^t+ e^{-t})^2}+ \frac{16}{(e^t+ e^{-t})^4}\right)
\end{align}



\end{document}